% Clase 27 --- Las tres formulas no son la misma (§1.3--1.4).
% A primer orden las tres coinciden; la diferencia recien se nota cuando la
% velocidad deja de ser chica frente a la de la onda.
\begin{center}
\begin{tikzpicture}

  \begin{axis}[
    fisfig, width=10.5cm, height=6.0cm,
    xlabel={$v_0/v$ \; (o $v_0/c$ en el caso de la luz)},
    ylabel={$f'/f$},
    xmin=0, xmax=0.62, ymin=0.95, ymax=2.9,
    xtick={0,0.2,0.4,0.6}, ytick={1,1.5,2,2.5},
    legend style={at={(0.03,0.97)}, anchor=north west},
  ]
    \addplot[figblue, smooth, samples=100, domain=0:0.6] {1+x};
    \addlegendentry{observador en movimiento: $1 + v_0/v$}
    \addplot[figred, smooth, samples=100, domain=0:0.6] {1/(1-x)};
    \addlegendentry{fuente en movimiento: $1/(1-v_0/v)$}
    \addplot[figamber, dashed, smooth, samples=100, domain=0:0.6]
      {sqrt((1+x)/(1-x))};
    \addlegendentry{luz (relativista): $\sqrt{(1+\beta)/(1-\beta)}$}

    \draw[figaux] (axis cs:0.12,0.95) -- (axis cs:0.12,2.9);
    \node[figlblaux, anchor=north west, align=left] at (axis cs:0.145,1.45)
      {a la izquierda de acá las tres\\ son indistinguibles};
  \end{axis}

\end{tikzpicture}\\[0.5em]
{\small\itshape Las dos fórmulas acústicas se parecen mucho pero no son la
misma: $1 + \beta$ contra $1/(1-\beta) = 1 + \beta + \beta^2 + \dots$, con
$\beta = v_0/v$. Difieren recién en el segundo orden, y por eso en la vida
cotidiana ---donde $v_0$ es una fracción chica de la velocidad del sonido--- da
lo mismo cuál se use. Que sean \textbf{distintas} en el fondo es lo importante:
significa que en las ondas mecánicas se puede saber, midiendo, quién se mueve
respecto del medio. Para la luz eso sería inaceptable, porque no hay medio y la
relatividad exige que sólo importe la velocidad \textbf{relativa}: la fórmula
correcta es la relativista, que es simétrica ---cambiarle el signo a $\beta$ la
invierte--- y que a primer orden se confunde con las otras dos,
$f' \simeq f(1 + v_0/c)$. Ése es el orden que alcanza para todas las
aplicaciones habituales: radar de velocidad, sonar y ecografía Doppler.}
\end{center}
